These are some common pitfalls users have reported when using Dyninst.
Many of these are either due to limitations in the current
implementations, or reflect design decisions that may not
produce the expected behavior from the system.

\setlength{\leftskip}{0.5cm}

  \textbf{Attach followed by detach}
  
  If a mutator attaches to a mutatee, and immediately exits, the
  current behavior is that the mutatee is left suspended.
  To make sure the application continues, call detach with the appropriate flags.
  
  \textbf{Attaching to a program that has already been modified by Dyninst}
  
  If a mutator attaches to a program that has already been modified by
  a previous mutator, a warning message will be
  issued. We are working to fix this problem, but the correct semantics
  are still being specified. Currently, a message
  is printed to indicate that this has been attempted, and the attach will fail.
  
  \textbf{Dyninst is event-driven}
  
  Dyninst must sometimes handle events that take place in the mutatee,
  for instance when a new shared library is loaded,
  or when the mutatee executes a fork or exec. Dyninst handles events
  when it checks the status of the mutatee, so to
  allow this the mutator should periodically call one of the functions
  \code{BPatch::pollForStatusChange}, \code{BPatch::wait­ForStatusChange},
  \code{BPatch\_thread::isStopped}, or \code{BPatch\_­thread::is­Termin­ated}.
  
  \textbf{Missing or out-of-date DbgHelp DLL (Windows)}
  
  Dyninst requires an up-to-date DbgHelp library on Windows. See the
  section on Windows-specific architectural
  issues for details.
  
  \textbf{Portland Compiler Group – missing debug symbols}
  
  The Portland Group compiler (pgcc) on Linux produces debug symbols
  that are not read correctly by Dyninst. The
  binaries produced by the compiler do not contain the source file
  information necessary for Dyninst to assign the debug
  symbols to the correct module.

\setlength{\leftskip}{0cm}
